\documentclass[10pt,a4paper]{amsart}
\usepackage[margin=19mm]{geometry}
\usepackage{amsmath,amssymb,mathtools}
\usepackage[scr=rsfs,cal=euler]{mathalfa}
\usepackage{xfrac}
\usepackage[unicode,hidelinks,bookmarks=false]{hyperref}
\newcommand{\ie}{i.e.\ }
\newcommand{\cf}{cf.\ }
\newcommand{\resp}{respectively\ }
\newcommand{\Subsection}{\subsection}
\providecommand{\nobreakdash}{\nobreak\hbox{-}}
\setlength{\emergencystretch}{2em}
\multlinegap0pt
\usepackage{amssymb}
\usepackage{bm}
\usepackage{mathtools}
\newtheorem{claim*}{Claim}
\DeclareMathOperator{\Hess}{Hess}
\newcommand{\M}{M}
\title{A Ballantine-type factorization of diffeomorphisms}
\author{Mahmoud Abdelgalil, Tryphon T. Georgiou}
\date{Source: arXiv:2609.06180v3 (CC BY 4.0)}
\begin{document}
\maketitle

\noindent\emph{Source and licence.} A Ballantine-type factorization of diffeomorphisms. Version arXiv:2609.06180v3 (CC BY 4.0), distributed by arXiv under the Creative Commons Attribution 4.0 International licence, http://creativecommons.org/licenses/by/4.0/, as recorded in the arXiv licence metadata for this exact version and shown on the abstract page https://arxiv.org/abs/2609.06180v3. Full licence notice: \texttt{licenses/arxiv-2609.06180-CC-BY-4.0.txt}.\par\smallskip

\noindent\textit{Editorial scope.} Counted unit: the article's claim* claim* (source0.tex lines 701--703). The article's complete original proof of it is delivered with it (source0.tex lines 705--748, one \texttt{proof} environment of 44 lines). No mathematics is written, repaired or re-derived: the statement and the proof are the article's own lines, in the article's own notation, and no retained line is substituted or re-flowed.  Retained uncounted as the source-local context the proof needs: lines 695--697.  Nothing was omitted from any retained span, so the package carries no omission record.  No retained line contains a citation, so the wrapper carries no bibliography and no bibliography entry.  No retained line contains a figure environment, an image-inclusion command or drawing code, so no image is delivered and the package declares no asset dependency.  Editorial formatting only, in addition: an AMS \texttt{amsart} preamble that restates the article's own declarations and macros used by the retained lines, namely the environments \texttt{claim*} on separate counters, together with the standard packages those lines need.  The retained environments print in this document as Claim* 1 in order of appearance, whereas the article prints the same items with the numbering of its own full document.  The complete native source of the article as distributed in the arXiv e-print for this exact version is bundled unchanged as \texttt{sources/209501/source0.tex}.
\begin{align}\label{eq:zero_order_span}
    \mathrm{span}_{\mathbb{R}}\{\xi^\sharp,(g^{-1}\Hess\psi^2)_x\xi^\sharp,\dots,(g^{-1}\Hess\psi^{m})_x\xi^\sharp\}=T_x\M,
\end{align}

\begin{claim*}
    The spanning condition \eqref{eq:zero_order_span} is satisfied by the family $\{\psi^2,\dots,\psi^m\}$.
\end{claim*}

\begin{proof}
    First note that, for all $i\in\{1,\dots,\tbinom{N+1}{2}\}$, the Leibniz rule applied to the Hessian gives
%
\begin{align*}
    g^{-1}\Hess \psi^{2N+i+1} = g^{-1}\Hess (\beta^k\beta^l) &= \beta^l\, g^{-1}\Hess \beta^k + \beta^k\, g^{-1}\Hess \beta^l + g^{-1}(d\beta^l\otimes d\beta^k + d\beta^k\otimes d\beta^l),
\end{align*}
%
where $(k,l)=\iota(i)$. The first two terms always belong to the pointwise linear span of $g^{-1}\Hess\psi^{i+1}$ for $i\in\{1,\dots,N\}$, and so it suffices to satisfy the span condition
%
\begin{align*}
    \mathbb{R}\,\xi^\sharp +  \mathrm{span}_{\mathbb{R}}\{g^{-1}(d\beta^l\otimes d\beta^k + d\beta^k\otimes d\beta^l)_x\xi^\sharp~|~1\leq k\leq l \leq N\} = T_x\M,
\end{align*}
%
for all $x\in\M$ and all $\xi\in T_x^*\M\backslash\{0\}.$ Since $(\beta^1,\dots,\beta^N)$ is an immersion, the differentials $\{d\beta^1,\dots,d\beta^N\}$ span the cotangent bundle pointwise everywhere and, therefore, the set of symmetric $(0,2)$-tensor fields 
$$
\{d\beta^l\otimes d\beta^k + d\beta^k\otimes d\beta^l ~|~1\leq k\leq l \leq N\}
$$
%
spans the set of all symmetric $(0,2)$-tensor fields on $\M$ pointwise, i.e., 
%
\begin{align*}
    \mathrm{span}_{\mathbb{R}}\{(d\beta^l\otimes d\beta^k + d\beta^k\otimes d\beta^l)_x~|~1\leq k\leq l \leq N\} = \mathrm{Sym}^2(T_x^*\M),
\end{align*}
%
for all $x\in \M$. For any $v\in T_x\M$, define the symmetric $(0,2)$-tensor $S\in\mathrm{Sym}^2(T_x^*\M)$ by
%
\begin{align*}
    S:= \xi\otimes v^\flat +v^\flat\otimes\xi,
\end{align*}
%
which, as discussed above, belongs to the pointwise span of the family. We proceed to compute that
%
\begin{align*}
    (g_x^{-1}S)\xi^\sharp = g_x(v,\xi^\sharp)\xi^\sharp+|\xi|_g^2v \implies v =  \frac{1}{|\xi|_g^2}((g_x^{-1}S)\xi^\sharp - g_x(v,\xi^\sharp)\xi^\sharp).
\end{align*}
%
which, since $\xi\neq 0$, is well-defined. Therefore, 
%
\begin{align*}
    v\in \mathbb{R}\,\xi^\sharp +  \mathrm{span}_{\mathbb{R}}\{g^{-1}(d\beta^l\otimes d\beta^k + d\beta^k\otimes d\beta^l)_x\xi^\sharp~|~1\leq k\leq l \leq N\},
\end{align*}
%
and, since $v\in T_x\M$ was arbitrary, the statement of the claim follows.
\end{proof} 

\end{document}
